\begin{frame}{자주 하는 실수 ①②: KL 항 계산}
  \begin{block}{① 복원 항 기대값에 $\sigma^2$ 항을 빼먹는 것}
    $\mathbb{E}[\log p(x|z)]$를 계산할 때 $(x-\mu)^2$만 쓰고 분산
    $\sigma_q^2$ 항을 빼먹으면 복원 항이 $-0.919$로 \textbf{과소평가}
    (정답 $-1.419$). $z = \mu + \sigma\epsilon$일 때 $\mathbb{E}[(x-z)^2]
    = (x-\mu)^2 + \sigma^2$ — ``평균으로 대체''한 것만으로 오차가
    사라지지 않음. $\sigma^2$가 큰(인코더가 불확실한) 입력에서 손실값이
    실제보다 크게 나타나는(= ELBO 과소평가) 방향으로 오류를 만들어,
    ``정규화 항을 줄이는데 왜 손실이 안 줄지''를 진단할 때 혼란의 원인.
  \end{block}
  \vskip0.4em
  \begin{block}{② KL 항을 ``분산만'' 계산하는 것}
    인코더가 $(\mu, \sigma^2)$를 함께 내놓는데, KL 항에 $\mu^2$(위치)
    항을 빼먹고 $\sigma^2 - 1 - \log\sigma^2$만 쓰면, \textbf{$\mu$에
    대한 그래디언트가 0} — 인코더의 평균 출력 위치에 대해 정규화 항이
    아무 작용도 하지 않아 ``분산만 조절된다''는, 학습이 이상하게
    멈추는 원인. $\frac{1}{2}(\sigma^2 + \mu^2 - 1 - \log\sigma^2)$의
    $\mu^2$ 항이 \textbf{인코더의 위치($\mu$)를 원점 방향으로
    끌어당기는 힘}임을 꼭 기억.
  \end{block}
\end{frame}
